\documentclass[11pt,reqno]{amsart}

\usepackage{amsmath,amsthm,amssymb}
\usepackage{mathtools}
\usepackage{hyperref}
\hypersetup{
  colorlinks=true,
  linkcolor=blue!50!black,
  citecolor=blue!50!black,
  urlcolor=blue!60!black
}
\usepackage{listings}
\usepackage{xcolor}
\usepackage{array}
\usepackage{booktabs}
\usepackage{geometry}
\geometry{margin=1in}

\lstset{
  basicstyle=\small\ttfamily,
  breaklines=true,
  columns=fullflexible,
  keepspaces=true,
  frame=single,
  framesep=4pt,
  rulecolor=\color{black!30},
  showstringspaces=false,
  aboveskip=6pt,
  belowskip=6pt
}

\newtheorem{theorem}{Theorem}[section]
\newtheorem{proposition}[theorem]{Proposition}
\newtheorem{lemma}[theorem]{Lemma}
\newtheorem{corollary}[theorem]{Corollary}
\theoremstyle{definition}
\newtheorem{definition}[theorem]{Definition}
\newtheorem{remark}[theorem]{Remark}

\newcommand{\CombLemma}{\texttt{CombLemma}}
\newcommand{\Mathlib}{\texttt{Mathlib}}
\newcommand{\unit}{\mathbf{I}}
\newcommand{\NC}{\mathrm{NC}}
\newcommand{\leanrepo}{https://github.com/MathNetwork/comb_arg/blob/v0.1.0/CombLemma}
% Inline, colored, clickable link to a verified Lean declaration.
% #1 = Lean declaration name (as rendered)
% #2 = file path inside CombLemma/ (e.g. SupReduction.lean)
\newcommand{\defref}[2]{%
  \href{\leanrepo/#2}{%
    \textcolor{green!45!black}{\ensuremath{\checkmark}}\hspace{0.15em}%
    \textcolor{blue!55!black}{\texttt{#1}}%
  }%
}

\title{Compiling the Almgren--Pitts Combinatorial Lemma}

\author{Xinze Li}
\address{\textit{Correspondence:} \texttt{lixinze@math.utoronto.ca}}
\date{\today}

\begin{document}

\begin{abstract}
The Almgren--Pitts combinatorial lemma is a recurring technical
step in the min-max theory of minimal hypersurfaces. We extract it
from the literature as a standalone metric-combinatorial algorithm
and formalize it in Lean 4 as an importable library.
\end{abstract}

\maketitle

\section{Introduction}
\label{sec:intro}

A minimal hypersurface is a critical point of the area
functional in codimension one. On a closed Riemannian manifold
$(M^{n+1}, g)$, existence of a nontrivial minimal hypersurface
is nontrivial: direct minimization over a fixed homology class
may produce only degenerate configurations, and the relevant
critical points are typically of saddle type rather than minima.
Birkhoff~(1917) established the prototype of a \emph{min-max}
construction by producing closed geodesics on the sphere via a
sweepout by loops; Almgren and Pitts, in the 1960s and 1970s,
generalized the construction to higher-dimensional varifolds.
The output is the following existence theorem, with regularity
supplied by Schoen--Simon~\cite{SS81}.

\begin{theorem}[Almgren--Pitts min-max theorem~\cite{Pit81,SS81}]
\label{thm:minmax}
Every closed Riemannian manifold $(M^{n+1}, g)$ admits a
nontrivial embedded minimal hypersurface $\Sigma \subset M$
whose singular set $\mathrm{Sing}\,\Sigma$ has Hausdorff
dimension at most $n-7$.
\end{theorem}

Theorem~\ref{thm:minmax} was established in its systematic form
by Pitts in the monograph~\cite{Pit81}; it has subsequently been
re-proved and extended in various frameworks by
Colding--De Lellis~\cite{CDL03}, De Lellis--Tasnady~\cite{DLT13},
Marques--Neves~\cite{MN14}, and others.

Pitts's proof splits into two parts. The first part executes the
min-max construction: starting from a sweepout family $\Lambda$
and the min-max value
\[
m_0(\Lambda) \;=\; \inf_{\Lambda} \, \max_{t} \, \mathcal{H}^n(\Gamma_t),
\]
one produces an almost-minimizing sequence of varifolds. The
second part is the regularity theory for the limiting varifold.
Within the first part, one key technical step is a combinatorial
assembly. In the formulation of De Lellis--Tasnady~\cite[\S3.2]{DLT13},
it reads as follows.

\begin{theorem}[De Lellis--Tasnady~{\cite[\S3.2]{DLT13}}]
\label{thm:comb-intro}
Let $\{\partial \Omega_t\}_{t \in [0,1]}$ be a sweepout of
$(M^{n+1}, g)$ with $\mathcal{H}^n(\partial \Omega_t)$
continuous in $t$, let $m_0 = \sup_t \mathcal{H}^n(\partial
\Omega_t)$, and let $N > 0$. Suppose that at each $t$ with
$\mathcal{H}^n(\partial \Omega_t) \ge m_0 - 1/N$ there is a
nested pair $A_t \subset\subset B_t \subset\subset U_t'$ in $M$
such that $\partial \Omega_t$ is almost minimizing in
$(A_t, B_t)$, and a replacement family
$\{\partial \Omega_{t,\tau}\}_{\tau \in [0,1]}$ (given
by~\cite[Lemma 3.1]{DLT13}) with $\partial \Omega_{t,0} =
\partial \Omega_t$ and $\mathcal{H}^n(\partial \Omega_{t,1})
\le \mathcal{H}^n(\partial \Omega_t) - 1/(4N)$ inside $U_t'$.
Then there exists a modified sweepout $\{\partial \Omega_t'\}$
with $\sup_t \mathcal{H}^n(\partial \Omega_t') \le m_0 - 1/(4N)$.
\end{theorem}

The combinatorial content of Theorem~\ref{thm:comb-intro} --- a
one-parameter assembly turning local near-critical savings into a
global supremum improvement --- recurs in Pitts's original proof
of Theorem~\ref{thm:minmax}~\cite[proof of Thm.~4.10]{Pit81} and
in the subsequent treatments~\cite[\S5]{CDL03},
\cite[Prop.~4.8]{DLR17}, \cite[Prop.~5.8]{DL12}, and
\cite[\S4]{MN14}, each time formulated in the geometric
language of the paper at hand.

The recurrence is structural. The statement of
Theorem~\ref{thm:comb-intro} is paper-agnostic: it depends on no
geometric measure theory, no admissible class, no ambient
regularity. Its content is a bookkeeping argument on continuous
real-valued functions, free of the geometric material that
varies between papers. Yet the statement has no canonical
reference, so each author re-derives it in local notation.

This note takes a different approach. We extract
Theorem~\ref{thm:comb-intro} as an independent statement and
formalize it in Lean 4 as the library \CombLemma. Informally,
the extracted form is the following.

\begin{theorem}[Extracted combinatorial lemma; formalized as \defref{exists\_sup\_reduction}{SupReduction.lean}]
\label{thm:comb-extracted}
Let $X$ be a pseudo-metric space equipped with an abstract
paired-cover structure. Let $f : [0,1] \to \mathbb{R}$ be
continuous with $m_0 = \sup f$ and $m_0 > 0$, and let $N > 0$.
Assume that at every $t \in [0,1]$ with $f(t) \ge m_0 - 1/N$
there is a \emph{local witness}: an open neighborhood $U_t \ni t$,
a paired cover $c_t$ in $X$, and a continuous replacement energy
$E_t : [0,1] \to \mathbb{R}$ satisfying $f(s) - E_t(s) \ge 1/(4N)$
for all $s \in U_t$. Then there exists $f' : [0,1] \to \mathbb{R}$
with $\sup f' \le m_0 - 1/(4N)$.
\end{theorem}

Theorem~\ref{thm:comb-extracted} is the scalar skeleton of
Theorem~\ref{thm:comb-intro}. The geometric content has been
stripped: the boundary $\partial \Omega_t$ is replaced by a
continuous function $f : [0,1] \to \mathbb{R}$; the
Hausdorff-measure hypothesis becomes $m_0 = \sup f$; the nested
pair $(A_t, B_t, U_t')$ collapses to a single open neighborhood
$U_t$; the replacement family
$\{\partial \Omega_{t,\tau}\}_{\tau \in [0,1]}$ collapses to a
single continuous $E_t$; the almost-minimizing assumption on
the varifold disappears. What remains is the combinatorial
bookkeeping: a near-critical cover, uniform saving bounds, and
the assembly of a globally improved $f'$. The abstract
paired-cover class on $X$ is a scaffolding parameter that
would, in a GMT instantiation, carry the geometric data of the
replacement; it is not load-bearing in the present proof (see
\S\ref{sec:discussion}). The formalization is machine-verified
against \Mathlib; it declares no custom axioms and contains no
\texttt{sorry}. A subsequent min-max paper that establishes its
own local-replacement lemma can then invoke the combinatorial
step as a one-line import, with the verification following
along.

\S\ref{sec:def-thm} gives the definitions and the main theorem;
\S\ref{sec:proof} sketches the proof; \S\ref{sec:discussion}
records the structural details the formalization makes explicit.

\section{Definitions and theorem}
\label{sec:def-thm}

Throughout this section, let $X$ be a pseudo-metric space carrying
an abstract paired-cover structure (Definition~\ref{def:pcover}).
Let $K$ be a topological space. The main theorem specializes to
$K = \unit$ (see Remark~\ref{rem:K-unit}).

\subsection{Paired covers and local witnesses}

\begin{definition}[Pairable cover class]
\label{def:pcover}
A \emph{pairable-cover} structure on $X$ consists of: a type
$\mathrm{Cover}$, two functions
$\mathrm{leftRegion}, \mathrm{rightRegion} : \mathrm{Cover}
\to 2^X$, a diameter function
$\mathrm{diameter} : \mathrm{Cover} \to \mathbb{R}_{\ge 0}$,
a disjointness axiom
$\mathrm{leftRegion}(c) \cap \mathrm{rightRegion}(c) = \emptyset$,
and a diameter-nesting axiom stating that overlapping combined
regions of diameter-comparable covers are ordered by
inclusion \defref{PairableCover}{Witness.lean}.\footnote{See
Remark~\ref{rem:scaffolding} on the scaffolding status of this
class.}
\end{definition}

\begin{definition}[Local witness]
\label{def:witness}
A \emph{local witness} at $t : K$ with saving threshold
$\varepsilon \in \mathbb{R}$ is a tuple consisting of: an open
set $U \ni t$ in $K$, a cover $c : \mathrm{Cover}$, a continuous
function $E : K \to \mathbb{R}$, and the uniform inequality
$\forall s \in U,\; f(s) - E(s) \ge \varepsilon$, where
$f : K \to \mathbb{R}$ is the ambient energy \defref{LocalWitness}{Witness.lean}.
\end{definition}

\begin{definition}[Near-critical set]
\label{def:NC}
For continuous $f : K \to \mathbb{R}$, scalar $m_0$, and
$N \in \mathbb{N}$, the \emph{$1/N$-near-critical set} is
$\NC = \{t \in K : f(t) \ge m_0 - 1/N\}$ \defref{nearCritical}{Refinement/InitialCover.lean}.
\end{definition}

\begin{proposition}[Compactness of the near-critical set]
\label{prop:NC-compact}
If $K$ is compact and $f$ is continuous, $\NC$ is closed in $K$,
hence compact. If additionally $m_0 = \sup f$ and $N > 0$, $\NC$
is nonempty \defref{isCompact\_nearCritical}{Refinement/InitialCover.lean}, \defref{nearCritical\_nonempty}{Refinement/InitialCover.lean}.
\end{proposition}

\subsection{Intermediate structures}

Fix now $K = \unit$.

\begin{definition}[Initial cover]
\label{def:ic}
An \emph{initial cover} for $(f, m_0, N)$ is a finite family
$\{I_i\}_{i=1}^n$ of open intervals in $\unit$, each
$I_i = (c_i - r_i, c_i + r_i) \cap \unit$, together with a
witness center $w_i \in \NC$ and a local witness at $w_i$ with
saving $1/(4N)$, such that
\begin{enumerate}
\item[(a)] $c_i + r_i < c_{i+2} - r_{i+2}$ for all valid $i$
  (non-overlap at index gap 2);
\item[(b)] $I_i \subset U_{w_i}$ (the $i$-th interval is
  contained in the $i$-th witness neighborhood);
\item[(c)] $\NC \subset \bigcup_i I_i$.
\end{enumerate}
We write this structure
$\mathrm{InitialCover}(X, f, m_0, N)$ \defref{InitialCover}{Refinement/InitialCover.lean}.\footnote{Conditions (a) and (b) transcribe DLT~\cite[\S3.2]{DLT13}'s spacing and subset conditions verbatim; the pair $(c_i, w_i)$ replaces DLT's single $t_i$, see Finding~2 in \S\ref{sec:discussion}.}
\end{definition}

\begin{definition}[Partial refinement]
\label{def:pr}
Given an initial cover $\mathrm{ic}$ and a natural $\ell$, a
\emph{partial refinement of length $\ell$} consists of pieces
$J_1,\ldots,J_\ell \subseteq \unit$ and an index assignment
$\sigma : \{1,\ldots,\ell\} \to \{1,\ldots,n\}$ satisfying
\begin{enumerate}
\item[(i)] $J_k \subseteq I_{\sigma(k)}$ for each $k$;
\item[(ii)] $I_{\sigma(k)} \subseteq \bigcup_{k'} J_{k'}$ for each
  $k$.
\end{enumerate}
We write this structure
$\mathrm{PartialRefinement}(\mathrm{ic}, \ell)$ \defref{PartialRefinement}{Refinement/PartialRefinement.lean}.
\end{definition}

\begin{definition}[Refinement]
\label{def:ref}
A \emph{refinement} of $(K, X, f, m_0, N)$ consists of a finite
nonempty index set $\iota$, pieces $\{p_l\}_{l \in \iota}$
covering $\NC$, a replacement-energy family
$\{E_l : K \to \mathbb{R}\}_{l \in \iota}$, a saving
family $\{s_l \in \mathbb{R}_{>0}\}_{l \in \iota}$, and three
invariants: (I) $f - E_l \ge s_l$ on $p_l$; (II) each point of
$K$ lies in at most two pieces; (III) $s_l \ge 1/(4N)$ for all
$l$ \defref{EnergyBound.Refinement}{EnergyBound.lean}.
\end{definition}

\subsection{Main theorem}

\begin{theorem}[$1/(4N)$ sup reduction]
\label{thm:main}
Let $X$ be a pseudo-metric space with a pairable-cover structure
(Definition~\ref{def:pcover}). Let $f : \unit \to \mathbb{R}$
be continuous, let $m_0 \in \mathbb{R}$ with $m_0 > 0$ and
$m_0 = \sup f$, and let $N \in \mathbb{N}$ with $N > 0$. Assume
that for every $t \in \NC$ there exists a local witness at $t$
with saving $1/(4N)$ (Definition~\ref{def:witness}). Then there
exists $f' : \unit \to \mathbb{R}$ with
\[
\sup f' \;\le\; m_0 - \frac{1}{4N}.
\]
\end{theorem}

\begin{proof}
Follows from Propositions~\ref{prop:exists-ic},
\ref{prop:exists-terminal},
and Theorem~\ref{thm:arithmetic} by
Lemma~\ref{lem:assembly}; see \S\ref{sec:proof}. The full chain
is \defref{exists\_sup\_reduction}{SupReduction.lean}, which
composes \defref{Refinement.exists\_refinement}{Refinement/Assembly.lean}
with \defref{EnergyBound.exists\_reducedEnergy\_sup\_lt}{EnergyBound.lean}.
\end{proof}

\begin{remark}
\label{rem:K-unit}
Theorem~\ref{thm:main} specializes to $K = \unit$. The
specialization enters Definitions~\ref{def:ic}~and~\ref{def:pr}
(intervals of $\unit$, $\sigma$ into a finite index set coming
from the cover of $\NC \subset \unit$). The arithmetic step
(Theorem~\ref{thm:arithmetic}) is formulated for general compact
$K$.
\end{remark}

\begin{remark}[Reduction, not contradiction]
\label{rem:reduction}
The theorem is stated as an existence of $f'$ rather than as a
contradiction $\bot$. This is deliberate: the hypothesis gives
$m_0 = \sup f$ for a specific $f$, not that $m_0$ is a
category-theoretic min-max level. Consumers who supply an
admissible class $\mathcal{A}$ with $m_0 = \inf_{g \in
\mathcal{A}} \sup g$ and $f \in \mathcal{A}$ layer the
contradiction on top by showing $f' \in \mathcal{A}$. See~\cite[\S3.2]{DLT13} for the GMT instance.
\end{remark}

\subsection{The arithmetic half}

The arithmetic half takes any refinement of $(K, X, f, m_0, N)$
and builds $f'$ by scalar subtraction.

\begin{theorem}[Arithmetic assembly]
\label{thm:arithmetic}
Let $K$ be compact and nonempty, $f : K \to \mathbb{R}$
continuous, $m_0 = \sup f$ with $m_0 > 0$, $N > 0$, and let
$R$ be a refinement of $(K, X, f, m_0, N)$
(Definition~\ref{def:ref}). Then there exists
$f' : K \to \mathbb{R}$ with $\sup f' \le m_0 - 1/(4N)$.
\end{theorem}

The proof constructs
$f'(t) := f(t) - \sum_{l \in \iota : t \in p_l} s_l$
and establishes the bound by a case split on whether $t \in \NC$,
using invariants~(II) and~(III) from
Definition~\ref{def:ref}. The formal
statement is \defref{EnergyBound.exists\_reducedEnergy\_sup\_lt}{EnergyBound.lean}.

\section{Proof structure}
\label{sec:proof}

We split the proof of Theorem~\ref{thm:main} into three stages.
Stage~A builds an initial cover from the witness hypothesis.
Stage~B runs an induction producing a terminal partial refinement
with injective $\sigma$. Stage~C assembles the refinement by
extending saving bounds to closures and deriving TwoFold.
Theorem~\ref{thm:arithmetic} closes.

\subsection{Stage A: existence of the initial cover}

\begin{proposition}[Existence of an initial cover]
\label{prop:exists-ic}
Under the hypotheses of Theorem~\ref{thm:main}, an initial cover
(Definition~\ref{def:ic}) exists \defref{exists\_initialCover}{Refinement/CoverConstruction.lean}.
\end{proposition}

\begin{proof}[Construction]
Compactness of $\NC$ (Proposition~\ref{prop:NC-compact}) plus the
open cover $\{U_t\}_{t \in \NC}$ afforded by the witness
hypothesis admit a Lebesgue number $\lambda > 0$: every
$\lambda$-ball in $\NC$ fits inside some $U_t$. Choose
$M \in \mathbb{N}$ with $M > 4/\lambda$, form the grid
$c_k = k/M$ for $k = 0, \ldots, M$, and select those $c_k$ within
$\lambda/4$ of $\NC$. For each selected $c_k$ pick a witness
center $w_k \in \NC$ within $\lambda/4$ of $c_k$; the ball of
radius $\lambda$ around $w_k$ fits inside some $U_{w_k}$, hence
so does the interval
$I_k := (c_k - 2/(3M), c_k + 2/(3M)) \cap \unit$. Spacing
condition~(a) holds because consecutive selected grid indices
differ by at least~2 in the selected finset, so their centers
differ by at least $2/M > 4/(3M) = 2 \cdot (2/(3M))$. Coverage
condition~(c) holds because every $t \in \NC$ is within
$1/(2M) < 2/(3M)$ of its rounded grid neighbor, which is
selected. Condition~(b) holds by the Lebesgue argument.
\end{proof}

\subsection{Stage B: the refinement induction}

Starting from an initial cover, one builds pieces inductively.

\begin{definition}[Base case]
\label{def:step0}
$\mathrm{step\_zero}(\mathrm{ic})$ is the length-$1$ partial
refinement with
$J_1 := I_1$ and $\sigma(1) := 1$ \defref{step\_zero}{Refinement/PartialRefinement.lean}.
\end{definition}

\begin{definition}[Inductive step]
\label{def:step-succ}
Given a partial refinement $\mathrm{pr}$ of length $\ell+1$ and
an index $i^* \in \{1,\ldots,n\}$ with
$I_{i^*} \not\subseteq \bigcup_k J_k$, define a length-$(\ell+2)$
extension by case split on whether
$I_{i^*} \subseteq J_\ell \cup I_{\sigma(\ell)}$:

\begin{itemize}
\item[(Case 1)] If yes, set $J_{\ell+1} := I_{i^*}$.
\item[(Case 2)] Otherwise, set
  $J_{\ell+1} := I_{i^*} \setminus I_{\sigma(\ell)}$.
\end{itemize}

In both cases set $\sigma(\ell+1) := i^*$. Both invariants of
Definition~\ref{def:pr} are preserved; in Case~2 the preservation
of (ii) uses the induction hypothesis on
$I_{\sigma(\ell)} \subseteq \bigcup_k J_k$ \defref{step\_succ\_at}{Refinement/Induction.lean}.\footnote{A variant \texttt{step\_succ} chooses $i^*$ internally via $\mathrm{Finset.min'}$; \texttt{step\_succ\_at} takes $i^*$ as an explicit argument and returns the extension bundled with the ``$i^*$ is covered'' property, for termination.}
\end{definition}

\begin{proposition}[Termination]
\label{prop:exists-terminal}
Starting from $\mathrm{step\_zero}(\mathrm{ic})$ and iterating
$\mathrm{step\_succ\_at}$, one reaches in at most $n$ steps a
partial refinement of some length $L$ with $\sigma$ injective
and $I_i \subseteq \bigcup_k J_k$ for every
$i$ \defref{exists\_terminal\_refinement}{Refinement/Induction.lean}.
\end{proposition}

\begin{proof}[Argument]
Let $\mathrm{rem}(\mathrm{pr})$ be the finite set of indices
$i$ with $I_i \not\subseteq \bigcup_k J_k$. One step of
$\mathrm{step\_succ\_at}$ at $i^* \in \mathrm{rem}(\mathrm{pr})$
(chosen as the smallest such) produces $\mathrm{pr}'$ with
$i^* \notin \mathrm{rem}(\mathrm{pr}')$ and
$\mathrm{rem}(\mathrm{pr}') \subseteq \mathrm{rem}(\mathrm{pr})
\setminus \{i^*\}$, hence strict decrease. Injectivity of
$\sigma$ follows from the case split in
Definition~\ref{def:step-succ}: the new $\sigma(\ell+1) = i^*$
is fresh because $i^* \in \mathrm{rem}$ while existing
$\sigma(k)$ correspond to already-covered indices.
\end{proof}

\subsection{Stage C: assembly}

The final stage extends the saving bound from open $J_k$ to
$\overline{J_k}$ and derives TwoFold from spacing~(a).

\begin{proposition}[TwoFold via parity]
\label{prop:twofold}
For a partial refinement $\mathrm{pr}$ of length $L$ with
$\sigma$ injective, every $t \in \unit$ lies in at most two of
$\overline{J_1}, \ldots, \overline{J_L}$ \defref{terminal\_twoFold}{Refinement/Assembly.lean}.
\end{proposition}

\begin{proof}[Argument]
From spacing condition~(a), iterated $m$ times, one obtains
that indices $i, j$ with $|i - j| = 2(m+1)$ (an even gap $\ge
2$) have strict separation between $\overline{I_i}$ and
$\overline{I_j}$, hence
$\overline{I_i} \cap \overline{I_j} = \emptyset$ \defref{InitialCover.chain\_spacing}{Refinement/Disjointness.lean}, \defref{InitialCover.closure\_disjoint\_of\_even\_gap}{Refinement/Disjointness.lean}.

Suppose three distinct indices $k_1, k_2, k_3 \in \{1,\ldots,L\}$
admit $t \in \overline{J_{k_1}} \cap \overline{J_{k_2}} \cap
\overline{J_{k_3}}$. Since $\overline{J_k} \subseteq
\overline{I_{\sigma(k)}}$ and $\sigma$ is injective,
$\sigma(k_1), \sigma(k_2), \sigma(k_3)$ are three distinct
indices in $\{1,\ldots,n\}$ with a common point in the
closures of $I_{\sigma(k_i)}$. Among any three distinct
integers $a < b < c$ some pair has even gap $\ge 2$: either
$c - a$ is even (and necessarily $\ge 2$), or $c - a$ is odd
and $\ge 3$, in which case $b - a$ and $c - b$ have opposite
parities and one of them is the even gap. The pair with even
gap has disjoint closures, contradicting the common
point \defref{InitialCover.not\_three\_overlap}{Refinement/Disjointness.lean}.
\end{proof}

\begin{proposition}[Closure extension of saving bound]
\label{prop:saving-bound}
Let $\mathrm{pr}$ be a partial refinement. For every $k$ and
every $t \in \overline{J_k}$,
$f(t) - E_{\sigma(k)}(t) \ge 1/(4N)$, where
$E_{\sigma(k)}$ is the replacement energy of the $\sigma(k)$-th
local witness \defref{saving\_bound\_closure}{Refinement/Assembly.lean}.
\end{proposition}

\begin{proof}[Argument]
On $J_k \subseteq U_{w_{\sigma(k)}}$ the witness gives
$f - E_{\sigma(k)} \ge 1/(4N)$. The function
$g := f - E_{\sigma(k)}$ is continuous
(Definitions~\ref{def:witness} requires $E_t$ continuous;
$f$ is continuous by hypothesis). A closure point is a limit of
a sequence in $J_k$, and $g$ preserves $\ge$-inequalities under
limits.
\end{proof}

\begin{lemma}[Assembly]
\label{lem:assembly}
A partial refinement $\mathrm{pr}$ of length $L$ with $\sigma$
injective and $I_i \subseteq \bigcup_k J_k$ for every $i$ yields
a refinement of $(\unit, X, f, m_0, N)$ in the sense of
Definition~\ref{def:ref}, with
$p_k := \overline{J_k}$, uniform saving
$s_k := 1/(4N)$, and
$E_k := $ the replacement energy of the
$\sigma(k)$-th witness \defref{Refinement.exists\_refinement}{Refinement/Assembly.lean}.
\end{lemma}

\begin{proof}[Argument]
Invariant~(I) is Proposition~\ref{prop:saving-bound}.
Invariant~(II) is Proposition~\ref{prop:twofold}.
Invariant~(III) holds by construction ($s_k = 1/(4N)$).
Coverage of $\NC$ chains $\NC \subseteq \bigcup_i I_i
\subseteq \bigcup_k J_k \subseteq \bigcup_k
\overline{J_k}$.
\end{proof}

Theorem~\ref{thm:main} now follows: chain
Propositions~\ref{prop:exists-ic}~and~\ref{prop:exists-terminal},
invoke Lemma~\ref{lem:assembly}, and apply
Theorem~\ref{thm:arithmetic}.

\section{Discussion}
\label{sec:discussion}

\subsection{Eight structural points}

Formalizing the proof surfaces eight structural details that the
prose exposition in~\cite{Pit81,CDL03,DLT13,DLR17,DL12,MN14}
compresses. None is a correction; each is a place where the
prose-level argument is locally underdetermined in a way the
formal proof must resolve.

\begin{enumerate}
\item The witness saving must commit to $\varepsilon = 1/(4N)$
  exactly, not $\exists \varepsilon > 0$. The weaker hypothesis
  does not thread through the lower bound on the
  sum-of-savings in Theorem~\ref{thm:arithmetic}.

\item The single $t_i$ of DLT~\cite[\S3.2]{DLT13} must be
  decoupled into an interval center $c_i$ and a witness
  center $w_i$ (Definition~\ref{def:ic}). The existence
  proof with $c_i = w_i$ forced is intractable under the
  standard Lebesgue-number approach: the grid point
  $c_k$ supplying the interval is not in general near-critical.

\item The index assignment $\sigma$ is not monotone
  (Definition~\ref{def:pr}, Proposition~\ref{prop:exists-terminal}).
  Intuition from geometric pictures suggests it should be; the
  smallest-index-selection rule occasionally backtracks when an
  earlier index has been covered as a side effect of a later
  step.

\item Spacing condition~(a) yields disjointness of
  $\overline{I_i}$ and $\overline{I_j}$ only for even gaps
  $|i - j| = 2(m+1)$. A counterexample with non-monotone
  centers (permitted by the abstract structure) witnesses
  overlap at odd gap~3.

\item TwoFold (Proposition~\ref{prop:twofold}) therefore
  requires an explicit parity argument: among three distinct
  indices some pair has even gap. DLT~\cite[\S3.2]{DLT13}
  asserts TwoFold in one sentence; the unpacking into
  chain-disjointness plus the parity rescue occupies a full
  proposition.

\item The \texttt{PairableCover} class's \texttt{diameter\_nesting}
  axiom is not invoked in the present proof. The scalar
  specialization to $\unit$ uses explicit radii instead. See
  Remark~\ref{rem:scaffolding}.

\item The replacement energy of a local witness must carry
  continuity explicitly (Definition~\ref{def:witness}). In GMT
  use this is automatic; under the abstraction, continuity is
  load-bearing in Proposition~\ref{prop:saving-bound} and must
  be an explicit field of \texttt{LocalWitness}.

\item The ambient class $X$ and \texttt{PairableCover} are
  carried in signatures but not consumed by the proof chain;
  see Remark~\ref{rem:scaffolding}.
\end{enumerate}

\subsection{Scaffolding}

\begin{remark}[\texttt{PairableCover} as scaffolding]
\label{rem:scaffolding}
The \texttt{PairableCover} class is carried through
\texttt{LocalWitness} and the intermediate structures but its
axioms are not used in the proof of Theorem~\ref{thm:main}.
Three outcomes are under consideration for v0.2: retain the
class as an extension point (future abstract-$K$ formulations
may consume it); rework the proof so the class is load-bearing
(following DLT's nested pair structure $A \subset\subset B
\subset\subset U'$ with annular shell $C = B \setminus A$);
or remove it in favor of a scalar-only interface. The v0.1
release defers the decision.
\end{remark}

\subsection{What is and is not in the repository}

\CombLemma~v0.1.0 contains the mathematics described above and
nothing else. In particular:
\begin{itemize}
\item It does \emph{not} construct local witnesses. The
  hypothesis of Theorem~\ref{thm:main} is supplied by the
  consumer.
\item It does \emph{not} contain any GMT content. Varifolds,
  Caccioppoli sets, currents, measure-theoretic regularity
  theorems are all absent.
\item It does \emph{not} prove the full min-max theorem. The
  local-replacement lemma and the min-max closure remain the
  consumer's work.
\item It does \emph{not} address multi-parameter sweepouts.
  Generalization to $K$ other than $\unit$ is future work;
  see \S\ref{ss:coverage} below.
\end{itemize}

\subsection{Coverage across the min-max literature}
\label{ss:coverage}

The formalization of \CombLemma~v0.1 specializes the sweepout
parameter space to $K = [0,1]$. Every min-max construction of
minimal hypersurfaces that operates on a one-parameter sweepout
consumes the same combinatorial assembly, and its scalar content
is what Theorem~\ref{thm:comb-extracted} formalizes. The list of
one-parameter constructions that \CombLemma~covers (as a
citation replacing the paper's internal re-derivation of the
combinatorial step) is summarized in Table~\ref{tab:coverage}.

\begin{table}[h]
\centering
\small
\begin{tabular}{@{}p{0.38\textwidth}p{0.56\textwidth}@{}}
\toprule
\textbf{Source} & \textbf{Where the one-parameter sweepout appears} \\
\midrule
Pitts~\cite[proof of Thm.~4.10]{Pit81}, $m=1$ case
  & The $(m,\mathbf{M})$-homotopy class of the main existence theorem, specialized to $m=1$. \\
Simon--Smith~\cite{Smi83}; Colding--De Lellis~\cite[\S5]{CDL03}
  & 1-parameter sweepout of surfaces in 3-manifolds, with continuous refinement. \\
De Lellis--Tasnady~\cite[\S3.2]{DLT13}
  & ``Standard notion of 1-parameter family of hypersurfaces,'' $\{\partial\Omega_t\}_{t\in[0,1]}$. \\
De Lellis--Ramic~\cite[Prop.~4.8]{DLR17}
  & 1-parameter sweepout with boundary. \\
Li--Zhou~\cite{LZ17}
  & Free-boundary 1-parameter sweepout $\phi : [0,1] \to \mathcal{Z}_n(\Omega,\partial\Omega)$. \\
Guaraco~\cite{Gua18}
  & 1-parameter Allen--Cahn min-max, $\Phi : [0,1] \to W^{1,2}(M)$. \\
Ketover~\cite{Ket16}
  & Equivariant 1-parameter Simon--Smith min-max. \\
Chambers--Liokumovich~\cite{CL16}
  & 1-parameter family of hypersurfaces for existence in complete, finite-volume manifolds. \\
Chambers--Liokumovich~\cite{CL20}
  & Monotonicity framework, $\{\partial\Omega_t\}_{t\in[0,1]}$. \\
Chodosh--Liokumovich--Spolaor~\cite{CLS22}
  & Non-excessive sweepout framework, explicitly ``one-parameter min-max.'' \\
Zhou~\cite{Zho19}
  & Multiplicity-one conjecture for 1-parameter sweepouts. \\
\bottomrule
\end{tabular}
\caption{One-parameter min-max works whose combinatorial step
is covered by \CombLemma~v0.1. In each, the paper's
local-replacement lemma supplies the hypothesis of
Theorem~\ref{thm:comb-extracted}; the combinatorial assembly
becomes a one-line invocation of
\texttt{CombLemma.exists\_sup\_reduction}.}
\label{tab:coverage}
\end{table}

Multi-parameter extensions are \emph{not} covered by v0.1. These
include Pitts's general $(m,\mathbf{M})$
framework~\cite[Thm.~4.10]{Pit81} with $m \ge 2$, the 5-parameter
Willmore construction of Marques--Neves~\cite{MN14}, the Weyl-law
and density work of Liokumovich--Marques--Neves~\cite{LMN18} and
Irie--Marques--Neves~\cite{IMN18}, and
De Lellis's multi-parameter exposition~\cite[Prop.~5.8]{DL12}.
A multi-parameter generalization is planned; see
\S\ref{ss:roadmap} for the specific interface changes it entails.

\subsection{Library roadmap}
\label{ss:roadmap}

Because \CombLemma's deliverable is a program and its public
interface (the \texttt{LocalWitness} structure, the statement of
\texttt{exists\_sup\_reduction}) is what downstream papers
consume, a ``future work'' list for this paper is a versioning
roadmap rather than a conjecture list. The planned trajectory is
as follows.

\begin{itemize}
\item \textbf{v0.2 (generalize $K$).} Move the parameter space
  from \texttt{unitInterval} to an arbitrary compact metric
  space. The \texttt{LocalWitness} structure is unchanged. The
  \texttt{InitialCover} structure is redesigned: spacing
  condition~(a), which currently relies on the linear order of
  $[0,1]$ (indices~$i$ and $i+2$), is replaced by a metric
  ball-covering condition on $K$ (pairs of balls at distance
  $\ge \lambda$ in the metric). The parity-rescue argument of
  \S\ref{sec:proof}~Stage~C becomes a disjointness argument on
  multiplicity in the metric cover. Downstream consumers who
  already operate on $[0,1]$ keep their existing call sites
  unchanged.

\item \textbf{v0.3 (multi-parameter $K = [0,1]^m$).} Move from
  arbitrary compact $K$ to a cubical-complex parametrization
  following Pitts~\cite{Pit81}. The \texttt{InitialCover} family
  of open intervals $\{I_i\}$ becomes a family of open cubes
  indexed by a cubical complex $I(m, n)$; the spacing condition
  becomes a face-relation bound; the parity argument becomes a
  cube-colouring argument counting cells by parity of
  coordinate-sum. \texttt{LocalWitness} and the main-theorem
  signature remain unchanged up to substituting~$K$; existing
  downstream consumers continue to compile.

\item \textbf{\texttt{PairableCover} decision.} The class is
  scaffolding in v0.1 (Finding~8, \S\ref{ss:points};
  Remark~\ref{rem:scaffolding}). v0.2 either (a) activates it by
  reworking Stage~C~(\S\ref{sec:proof}) to consume the
  diameter-nesting axiom in a DLT-style annular shell argument,
  or (b) removes it in favour of a scalar-only interface.
  Selection between~(a) and~(b) will be driven by whichever form
  a concrete downstream consumer establishes first.

\item \textbf{\texttt{Mathlib} integration.} \CombLemma~is
  currently a standalone project with its own package name.
  If the material is suitable for upstreaming, a future version
  will adapt naming to Mathlib conventions (namespace layout,
  camel/snake case per declaration type, split granularity
  matching Mathlib's file-per-concept norm) and submit the
  library as a \Mathlib~pull request. The adaptation will not
  change the public API names; it will re-export them under
  Mathlib's namespace structure.
\end{itemize}

None of the above is a mathematical open problem. Each is a
defined interface change with a concrete scope, a target version,
and an identified downstream compatibility story.

\begin{thebibliography}{Mathlib}

\bibitem[Pit81]{Pit81}
J.~T. Pitts, \emph{Existence and regularity of minimal surfaces
on Riemannian manifolds}, Mathematical Notes~27, Princeton
University Press, 1981.

\bibitem[SS81]{SS81}
R.~Schoen and L.~Simon, \emph{Regularity of stable minimal
hypersurfaces}, Comm.\ Pure Appl.\ Math.~34 (1981), no.~6,
741--797.

\bibitem[CDL03]{CDL03}
T.~H. Colding and C.~De Lellis, \emph{The min-max construction
of minimal surfaces}, Surv.\ Differ.\ Geom.~VIII (2003),
75--107.

\bibitem[DLT13]{DLT13}
C.~De Lellis and D.~Tasnady, \emph{The existence of embedded
minimal hypersurfaces}, J.\ Differential Geom.~95 (2013),
no.~3, 355--388.

\bibitem[DLR18]{DLR17}
C.~De Lellis and J.~Ramic, \emph{Min-max theory for minimal
hypersurfaces with boundary}, Ann.\ Inst.\ Fourier~68 (2018),
no.~5, 1909--1986.

\bibitem[DL16]{DL12}
C.~De Lellis, exposition of the combinatorial step in the min-max
construction, Ann.\ Inst.\ Fourier, 2016. (Precise article
reference to be confirmed against the author's library.)

\bibitem[MN14]{MN14}
F.~Marques and A.~Neves, \emph{Min-max theory and the Willmore
conjecture}, Ann.\ of Math.~(2) 179 (2014), no.~2, 683--782.

\bibitem[Smi83]{Smi83}
F.~Smith, \emph{On the existence of embedded minimal 2-spheres
in the 3-sphere, endowed with an arbitrary Riemannian metric},
PhD thesis, University of Melbourne, 1983.

\bibitem[LZ17]{LZ17}
M.~M.-C. Li and X.~Zhou, \emph{Min-max theory for free boundary
minimal hypersurfaces I: Regularity theory}, J.\ Differential
Geom.~118 (2021), no.~3, 487--553.

\bibitem[Gua18]{Gua18}
M.~A.~M. Guaraco, \emph{Min-max for phase transitions and the
existence of embedded minimal hypersurfaces}, J.\ Differential
Geom.~108 (2018), no.~1, 91--133.

\bibitem[Ket16]{Ket16}
D.~Ketover, \emph{Equivariant min-max theory}, arXiv:1612.08692
(2016).

\bibitem[CL16]{CL16}
G.~R. Chambers and Y.~Liokumovich, \emph{Existence of minimal
hypersurfaces in complete manifolds of finite volume},
Invent.\ Math.~219 (2020), no.~1, 179--217.

\bibitem[CL20]{CL20}
G.~R. Chambers and Y.~Liokumovich, \emph{Monotonicity of sweepouts
and minimal hypersurfaces}, arXiv preprint (2020).

\bibitem[CLS22]{CLS22}
O.~Chodosh, Y.~Liokumovich, and L.~Spolaor, \emph{Non-excessive
sweepouts and the min-max minimal hypersurface problem},
arXiv preprint (2022).

\bibitem[Zho19]{Zho19}
X.~Zhou, \emph{On the multiplicity one conjecture in min-max
theory}, Ann.\ of Math.~(2) 192 (2020), no.~3, 767--820.

\bibitem[LMN18]{LMN18}
Y.~Liokumovich, F.~C. Marques, and A.~Neves, \emph{Weyl law for
the volume spectrum}, Ann.\ of Math.~(2) 187 (2018), no.~3,
933--961.

\bibitem[IMN18]{IMN18}
K.~Irie, F.~C. Marques, and A.~Neves, \emph{Density of minimal
hypersurfaces for generic metrics}, Ann.\ of Math.~(2) 187 (2018),
no.~3, 963--972.

\bibitem[dMU21]{Lean4}
L.~de~Moura and S.~Ullrich, \emph{The Lean~4 theorem prover and
programming language}, in: CADE-28, Lecture Notes in Computer
Science~12699, Springer, 2021, 625--635.

\bibitem[MC20]{Mathlib}
The \Mathlib~Community, \emph{The Lean Mathematical Library},
in: CPP 2020, 367--381.

\end{thebibliography}

\appendix

\section{Repository structure}
\label{app:structure}

\CombLemma~is a Lean 4 \texttt{lake} project with a pinned
\Mathlib~dependency (commit recorded in
\texttt{lake-manifest.json}). The source layout is:

\begin{lstlisting}
CombLemma.lean          top-level re-export
CombLemma/
  Witness.lean          PairableCover, LocalWitness
  EnergyBound.lean      EnergyBound.Refinement, arithmetic
  SupReduction.lean     exists_sup_reduction (main theorem)
  Refinement.lean       facade re-exporting six submodules
  Refinement/
    InitialCover.lean   nearCritical, InitialCover
    CoverConstruction.lean  exists_initialCover
    PartialRefinement.lean  PartialRefinement, step_zero
    Induction.lean      step_succ, step_succ_at, exists_terminal_refinement
    Disjointness.lean   chain_spacing, disjoint_of_even_gap,
                        closure_disjoint_of_even_gap, not_three_overlap
    Assembly.lean       terminal_twoFold, saving_bound_closure,
                        exists_refinement
test/Smoke.lean         trivial PairableCover on R
examples/MinimalUsage.lean  invocation pattern for consumers
\end{lstlisting}

Thirteen Lean 4 files, approximately 1937 lines including
docstrings and copyright headers; no individual file exceeds
400 lines.

\section{Build and axiom audit}
\label{app:axioms}

The library builds with
\begin{lstlisting}
lake exe cache get   # fetch Mathlib oleans (first run)
lake build           # build CombLemma
lake build test      # build smoke test
lake build examples  # build examples
\end{lstlisting}
No custom axioms are declared. The axiom audit
\begin{lstlisting}
#print axioms CombLemma.exists_sup_reduction
\end{lstlisting}
reports dependence on three axioms: \texttt{propext},
\texttt{Classical.choice}, \texttt{Quot.sound}. These are
Lean 4's foundational axioms; all three are used throughout
\Mathlib.

\section{GMT $\leftrightarrow$ abstract glossary}
\label{app:glossary}

\begin{center}
\small
\begin{tabular}{@{}p{0.38\textwidth}p{0.56\textwidth}@{}}
\toprule
\textbf{GMT / min-max term} & \textbf{\CombLemma~correspondence} \\
\midrule
Sweepout & Continuous $f : K \to \mathbb{R}$ \\
Sweepout parameter space & $K$ (specialized to $\unit$) \\
Width of the sweepout & $m_0 = \sup f$ \\
Near-critical slice & $t$ with $f(t) \ge m_0 - 1/N$ \\
Replacement family at $t$ & Local witness at $t$ (Definition~\ref{def:witness}) \\
Energy of the replacement & Replacement-energy function $E_t$ \\
Quantitative saving & Threshold $\varepsilon$; here $\varepsilon = 1/(4N)$ \\
Varifold / current / Caccioppoli set & Ambient $X$ (abstracted away) \\
Paired-cover structure on $X$ & \texttt{PairableCover}~$X$ \\
Improved sweepout & $f'$ (conclusion of Theorem~\ref{thm:main}) \\
Min-max contradiction & Downstream consumer responsibility \\
\bottomrule
\end{tabular}
\end{center}

\section{Installation and invocation}
\label{app:install}

To consume \CombLemma~in a downstream Lean 4 project, add to
the project's \texttt{lakefile}:
\begin{lstlisting}
require CombLemma from git
  "https://github.com/MathNetwork/comb_arg.git" @ "v0.1.0"
\end{lstlisting}
Run \texttt{lake exe cache get} and \texttt{lake build}. In
consumer code:
\begin{lstlisting}
import CombLemma

example
    {X : Type*} [PseudoMetricSpace X] [CombLemma.PairableCover X]
    {f : unitInterval -> R} (hf : Continuous f)
    {m0 : R} (hm_pos : 0 < m0) (hm : m0 = sSup (Set.range f))
    {N : N} (hN : 0 < N)
    (witness : forall t : unitInterval, f t >= m0 - 1/(N:R) ->
        Nonempty (CombLemma.LocalWitness unitInterval X f t
                    (1 / (4 * (N : R))))) :
    exists f' : unitInterval -> R,
        sSup (Set.range f') <= m0 - 1/(4 * (N : R)) :=
  CombLemma.exists_sup_reduction hf hm_pos hm hN witness
\end{lstlisting}
(ASCII substitutions for Unicode in the Lean 4 source:
\texttt{R}~$= \mathbb{R}$, \texttt{N}~$= \mathbb{N}$,
\texttt{->}~$= \to$, \texttt{>=}~$= \ge$, \texttt{<=}~$= \le$,
\texttt{in}~$= \in$.)

\end{document}
